\chapter*{Úvod}
\addcontentsline{toc}{chapter}{Úvod}
\label{chap:introduction}

Tato práce dokumentuje průběh individuální odborné praxe, kterou jsem absolvoval ve společnosti AVE~soft s.r.o. v rámci závěrečné fáze bakalářského studia na Fakultě elektrotechniky a informatiky Vysoké školy báňské -- Technické univerzity Ostrava. Volba absolvovat odbornou praxi namísto klasické bakalářské práce vycházela z mého zájmu ověřit si teoretické znalosti v reálném prostředí softwarové firmy a získat zkušenost s vývojem produkčního kódu v rámci rozsáhlého informačního systému.

\section*{Motivace a cíle praxe}

Hlavním technickým cílem praxe byla modernizace modulu \textbf{Importér} -- nástroje pro hromadný import dat do informačního systému Evolio. Úkolem bylo navrhnout a implementovat nové webové rozhraní (Single Page Application) v technologii Angular, přispět k vývoji backendové vrstvy v podobě C\# doplňku (Add-in) a podílet se na návrhu databázového schématu.

Vedle hlavního projektu byly součástí praxe dva další významné okruhy úkolů: implementace uživatelského rozhraní pro AI asistenta s podporou streamované komunikace a analýzy nahraných dokumentů, a tvorba automatizovaných UI testů v jazyce C\# pomocí frameworku Selenium. Společným jmenovatelem všech činností bylo zapojení do kompletního vývojového cyklu firmy, včetně code review, verzování, průběžné integrace a nasazování (CI/CD).

\section*{Struktura bakalářské práce}

Práce je členěna následovně. Kapitola \ref{chap:company_profile} představuje společnost AVE~soft s.r.o. a moji pracovní pozici v rámci vývojového týmu. Kapitola \ref{chap:assignment_schedule} specifikuje zadané úkoly a jejich časový rozvrh. Kapitola \ref{chap:importer_implementation} tvoří technické jádro práce -- popisuje realizaci modulu Importér ve všech třech vrstvách (databázové, aplikační a prezentační). Kapitola \ref{chap:fe_issues_ai} se věnuje vývoji AI asistenta a řešení vybraných chyb ve stávajícím frontendu. Kapitola \ref{chap:testing} dokumentuje metodiku a implementaci automatizovaných testů. Kapitoly \ref{chap:studies_connection} a \ref{chap:competencies} reflektují vazbu praxe na studium a shrnují nově osvojené znalosti i zjištěné nedostatky ve vzdělání. Kapitola \ref{chap:evaluation} hodnotí dosažené výsledky a Závěr shrnuje přínos praxe.
